% ======================================================================
% Zeilenabstand und Absatzlayout
% ======================================================================
\setstretch{1.5} % 1,5-zeiliger Zeilenabstand im Fließtext
\AtBeginDocument{%
  \setlength{\parindent}{0pt}% kein Absatzeinzug
  \setlength{\parskip}{4pt plus 1pt minus 1pt}% Abstand zwischen Absätzen (leicht flexibel)
}


% ======================================================================
% Seitenränder (geometry)
% ======================================================================
\usepackage{geometry}
\geometry{paper=a4paper,left=30mm,right=25mm,top=25mm,bottom=15mm}

% Alternative Rand-Einstellungen (nur bei Bedarf verwenden)
% \geometry{paper=a4paper,left=35mm,right=25mm,top=25mm,bottom=25mm}


% ======================================================================
% Kopf- und Fußzeilen (scrlayer-scrpage)
% ----------------------------------------------------------------------
% Zentraler Seitenstil: Kopfzeile mit Headmark und Seitenzahl,
% Fußzeile vollständig leer.
% ======================================================================
\clearpairofpagestyles

% Kopfzeile: innen/links Headmark (kursiv), außen/rechts Seitenzahl
\ihead[\textit{\headmark}]{\textit{\headmark}}
\ohead[\pagemark]{\pagemark}

% Fußzeile deaktivieren
\lefoot{}\cefoot{}\refoot{}
\lofoot{}\cofoot{}\rofoot{}

% Schrift im Kopf: normal (nicht smallcaps o. ä.)
\renewcommand{\headfont}{\normalfont}

% Abmessungen/Abstände der Kopf- und Fußbereiche
\setlength{\headheight}{8mm}
\setlength{\headsep}{4mm}
\setlength{\footskip}{8mm}

% Kapitelanfangsseiten (plain) an den Standard-Seitenstil anpassen
% Hinweis: Nach dem Laden von scrlayer-scrpage ausführen.
\makeatletter
\let\ps@plain\ps@scrheadings
\makeatother

% Kopfbreite exakt auf Textbreite setzen
\AtBeginDocument{\KOMAoptions{headwidth=text}}


% ======================================================================
% Kapitelstart: vertikaler Abstand (KOMA-Script)
% ----------------------------------------------------------------------
% Schiebt die Kapitelüberschrift näher an den oberen Rand.
% Vorsicht: Zu große negative Werte können Layout-Probleme verursachen.
% ======================================================================
\renewcommand*\chapterheadstartvskip{\vspace*{-1.0cm}}


% ======================================================================
% Typografie: Satzzeichenabstände
% ----------------------------------------------------------------------
% Aktiviert „französischen Abstand“: einheitliche Abstände nach Satzzeichen.
% ======================================================================
\frenchspacing


% ======================================================================
% Schusterjungen und Hurenkinder vermeiden (Witwen/Waisen)
% ======================================================================
\clubpenalty=10000
\widowpenalty=10000
\displaywidowpenalty=10000


% ======================================================================
% Quellcode-Ausgabe (listings): Grundlayout
% ======================================================================

% Beispiel-Optionen (bei Bedarf aktivieren/anpassen)
% \lstset{numbers=left, numberstyle=\tiny, numbersep=5pt, breaklines=true}
% \lstset{emph={square}, emphstyle=\color{red}, emph={[2]root,base}, emphstyle={[2]\color{blue}}}

% Hintergrundfarbe für Listings
\definecolor{gray}{rgb}{0.9,0.9,0.9}

% Standard-Listing-Einstellungen
\lstset{
  basicstyle=\small\ttfamily,language={[LaTeX]TeX},
  numbersep=5mm,
  numbers=left,
  numberstyle=\tiny,
  breaklines=true,
  framexleftmargin=8mm,
  xleftmargin=8mm,
  backgroundcolor=\color{gray},
  captionpos=b
}


% ======================================================================
% Fußnoten: fortlaufende Nummerierung über Kapitel hinweg
% ======================================================================
\counterwithout{footnote}{chapter}


% ======================================================================
% Theorem-Umgebungen
% ======================================================================
\newtheorem{definition}{Definition}


% ======================================================================
% Abstände um Überschriften (titlesec)
% ----------------------------------------------------------------------
% Format: \titlespacing*{<Befehl>}{<linker Einzug>}{<Abstand davor>}{<Abstand danach>}
% ======================================================================
\titlespacing*{\chapter}{0pt}{-5pt}{12pt}
\titlespacing*{\section}{0pt}{12pt}{6pt}
\titlespacing*{\subsection}{0pt}{10pt}{4pt}
\titlespacing*{\subsubsection}{0pt}{8pt}{4pt}

% Nummerlose Kapitel (z. B. Abstract) separat definieren
\titlespacing*{name=\chapter,numberless}{0pt}{2pt}{10pt}
